\scanpage{45}
% archive foliation: 44
\uline{Prop}

$A$ $\mathbb{Q}$-alg. comm. unif.

$u$ un autom de $\mathbb{Q}$-alg., $n\in\mathbb{N}$,

$J$ l'idéal engendré par les $u(x)-x$ $(x\in A)$
\margin{\uline{NB} $u$ induit l'identité sur $\operatorname{gr}_J(A)$}

1) Conditions équivalentes

(i) $u$ op. diff d'ordre $\leq n$

\margin{\deleted{(ii') $u$ \unsure{induit} l'identité}}
(ii) $J^{n+1}=0$

\margin{\deleted{(iii') $\exists\unsure{I\subset A}$ que $I^{n+1}=0$, et que $u$ \unsure{induise} l'identité sur $\operatorname{gr}_I(A)$ \unsure{(0)}}}
(iii) $\exists I$ idéal \inserted{de $A$} tel que $I^{n+1}=0$, et que $u$ induit l'identité dans $A/I$, i.e.
\[
\uline{$ux-x$}\in I\qquad\forall x\in A\qquad[\text{i.e. }I\supset J]
\]
\deleted{(iv)} \inserted{Pour ceci, il \uline{suffit} que $(u-\operatorname{id})^{n+1}=0$, ou (ce qui revient au $=$)}
$\exists$ dérivation $\Theta:A\longrightarrow A$ avec
\[
\Theta^{n+1}=0,\qquad u=\exp\Theta
\]
\margin{est-ce nécessaire ?}

\deleted{\uline{NB} Si $N(A)$ est nilpotent, \unsure{ces cond\textsuperscript{ts} st équivalentes|ces cond\textsuperscript{ts} sont équivalentes}}

2) L'application
\[
\Theta\longmapsto\exp\Theta
\]
est une bijection de l'ens des dérivations $\Theta$ de $A$ telles que $\Theta^{n+1}=0$, dans l'ens des autom $u$ de $A$ \inserted{\unsure{telles|tels} que $(u-\operatorname{id})^{n+1}=0$ (\unsure{donc}} satisfaisant les conditions équivalentes plus haut).

\uline{Cor} Conditions équivalentes sur autom $u$

(i) $u$ op. diff

(ii) $J$ idéal nilp.

(iii) $\exists I$ nilpotent tel que $u$ induise $\operatorname{id}_{A/I}$

(iv) $\exists\Theta$ dérivation, nilpotente, telle que $u=\exp\Theta$

De plus, $\Theta\longmapsto\exp\Theta$ est une bij de l'ens des dérivations nilpotentes de $A$, et l'ens des autom \inserted{de $A$} qui sont des op. diff.
